\documentclass[10pt,a4paper]{amsart}
\usepackage[margin=19mm]{geometry}
\usepackage{amsmath,amssymb,mathtools}
\usepackage[scr=rsfs,cal=euler]{mathalfa}
\usepackage{xfrac}
\usepackage[unicode,hidelinks,bookmarks=false]{hyperref}
\newcommand{\ie}{i.e.\ }
\newcommand{\cf}{cf.\ }
\newcommand{\resp}{respectively\ }
\newcommand{\Subsection}{\subsection}
\providecommand{\nobreakdash}{\nobreak\hbox{-}}
\setlength{\emergencystretch}{2em}
\multlinegap0pt
\usepackage{bm}
\usepackage{bbm}
\usepackage{dsfont}
\usepackage{xspace}
\usepackage{cancel}
\usepackage{mathtools}
\usepackage{amsthm}
\makeatletter
\@ifundefined{theorem}{\newtheorem{theorem}{Theorem}}{}
\@ifundefined{proposition}{\newtheorem{proposition}[theorem]{Proposition}}{}
\makeatother
\def\ws{\widetilde{\Sigma}}
\title[Liftable braids and the coloured braid groupoid]{Liftable braids and the coloured braid groupoid}
\author{Joan Licata, Vera V\'ertesi}
\date{Source: arXiv:2508.05146v1 (CC BY 4.0)}
\begin{document}
\maketitle

\noindent\emph{Source and licence.} Liftable braids and the coloured braid groupoid. Version arXiv:2508.05146v1 (CC BY 4.0), distributed under the Creative Commons Attribution 4.0 International licence, as recorded in the arXiv licence metadata for this exact version and shown on the abstract page https://arxiv.org/abs/2508.05146v1. Full rights notice: licenses/arxiv-2508.05146-CC-BY-4.0.txt.\par\smallskip

\noindent\textit{Editorial scope.} Counted unit: the Theorem whose source label is \texttt{thm:homom2} in the article (source0.tex lines 570--571, source label \texttt{thm:homom2}), delivered together with the article's own complete original proof of it (source0.tex lines 573--590, one \texttt{proof} environment of 18 lines). No mathematics is written, repaired or re-derived: statement and proof are the article's own text line for line. Required context retained uncounted: prop:braidrel (lines 436--437). Every definition, display and label of those lines is retained unchanged. Omitted from this package and not redistributed: 1--435, 438--569, 572--572, 591--782. Nothing inside a retained excerpt is omitted or altered: every retained body is delivered exactly as the article's lines are written, so no omission receipt of any kind accompanies this package. Citations: the retained excerpts carry no citation command at all, so no bibliography entry of the article is transcribed and no reference is redistributed. Reference adaptation: none was needed and none was made; every reference inside the delivered unit resolves inside it, so no unresolved reference or citation is produced. Printed slips kept literal: no printed word, symbol, hypothesis or number of the counted statement or of its proof was altered. Layout and class: the article's own class is \texttt{amsart} with packages including palatino, pinlabel, pdfpages, hyperref, multicol, tikz-cd, amsmath, amssymb, amsthm, verbatim, amsfonts, amscd, flafter, epsf; the wrapper is the lane's pinned \texttt{amsart} editorial wrapper at 10pt on A4, which restates the theorem-like environments the retained text needs and adds the article's own macro definitions that the retained text needs (\texttt{\textbackslash ws}), with extra wrapper packages (bm, bbm, dsfont, xspace, cancel, mathtools). The delivered numbering is the unit's own. Assets: no retained span contains a figure environment, a graphics-inclusion command, a PDF-page inclusion, a table or a TikZ picture; no image file is copied into this package, the package declares no asset dependency and no image file is redistributed with it. Licence: version arXiv:2508.05146v1 of this article is distributed under the Creative Commons Attribution 4.0 International licence, as recorded in the arXiv licence metadata for this exact version and shown on the abstract page https://arxiv.org/abs/2508.05146v1. A full rights notice is delivered at licenses/arxiv-2508.05146-CC-BY-4.0.txt. Characters: every retained excerpt is pure ASCII, so the delivered engine, XeLaTeX with its default Latin Modern text font, reports no missing character.
\begin{proposition}\label{prop:braidrel} Suppose that $\beta=\sigma_{i_m}\dots\sigma_{i_1}$ and $\beta=\sigma_{j_p}\dots\sigma_{j_1}$ are two factorisations of $\beta$.  Then $\widetilde{\sigma}_{i_m}\dots \widetilde{\sigma}_{i_1}=\widetilde{\sigma}_{j_p}\dots \widetilde{\sigma}_{j_1}$.
\end{proposition}

\begin{theorem}\label{thm:homom2} The map $\Phi: \text{ColB}_{n,d}\rightarrow \mathcal{M}(\mu)$ is a groupoid homomorphism.
\end{theorem}

\begin{proof} 
Suppose that $\sigma_{i_2}\sigma_{i_1}$ is a coloured braid acting on the disc with initial labels $\tau$.  Recall that $\Phi(\sigma_{i_1})$ is a homeomorphism from $\ws_\tau$ to $\ws_{\sigma_{i_1}\cdot \tau}$ and $\Phi(\sigma_{i_2})$ is a homeomorphism from $\ws_{\sigma_{i_1}\cdot \tau}$ to $\ws_{\sigma_{i_2}\sigma_{i_1}\cdot \tau}$. 

Since $\Phi(\beta)$ is independent of the factorisation of $\beta$ as a product of generators, it suffices to show that for any Artin generators or their inverses, \begin{equation}\label{eq}\Phi(\sigma_{i_2}\sigma_{i_1})=\Phi(\sigma_{i_2})\circ\Phi(\sigma_{i_1}).\end{equation}
 
Furthermore, we must check that a trivial coloured braid maps to a trivial mapping class in $\mathcal{M}(\mu)$.  However, this follows immediately from Proposition~\ref{prop:braidrel}, which proved that trivial braids lift to identity morphisms in $\mathcal{G}(\mu)$.  Since $\phi(\beta) \cdot O_\tau$ is isotopic to $O_\tau$ for any trivial $\beta$, we see that $\Phi(\beta)$ is the identity morphism in $\mathcal{M}$.



A standard way to compare two homeomorphisms of a surface is to see that they act the same way on a sufficiently complicated set of curves.  In the present situation, we claim that the left- and right-hand sides of Equation~\ref{eq} act the same way on the lifts of the cut curves $\Delta$.  As with $\pi^{-1}(\Gamma)$, we discard components of $\pi^{-1}(\Delta)$ that do not meet at branch points, and we denote the remaining set of properly embedded arcs by $\widetilde{\Delta}$.  

Each of the three homeomorphisms in Equation~\ref{eq} is defined as a map sending a neighbourhood of one rigid vertex graph to a neighbourhood of another.  On the left-hand side, $\widetilde{\Gamma}_\tau$ is first altered by the $\mathcal{G}$-morphism $\phi(\sigma_{i_2}\sigma_{i_1})$ and then the surface neighbourhood of this new rigid vertex graph is mapped to $\widetilde{\Gamma}_{\sigma_{i_2}\sigma_{i_1}\cdot \tau}$ by $\Phi(\sigma_{i_2}\sigma_{i_1})$.  The intersection pattern between $\widetilde{\Gamma}_{\sigma_{i_2}\sigma_{i_1}\cdot \tau}$ and $\widetilde{\Delta}$ on $\ws_\tau$  is dictated by the specific arcslides and index swaps of $\phi(\sigma_{i_2}\sigma_{i_1})$, and this intersection pattern determines $\Phi(\sigma_{i_2}\sigma_{i_1})(\widetilde{\Delta})$.  

The right-hand side of Equation~\ref{eq} is a composition of two homeomorphisms, but the key observation is that the same arcslides and index swaps are applied to the rigid vertex graphs as in the first case.  There is an intermediate step that identifies a neighbourhood of $\phi(\sigma_{i_1})\cdot \widetilde{\Gamma}_\tau$ with $\widetilde{\Gamma}_{\sigma_{i_1}\cdot \tau}$ on $\widetilde{\Sigma}_{\sigma_{i_1}\cdot \tau}$, but this doesn't affect how the set of arcslides and index swaps associated to $\phi(\sigma_{i_2})$ change the intersection pattern with $\widetilde{\Delta}$.  Since  $\widetilde{\Gamma}_{\sigma_{i_2}\sigma_{i_1}\cdot \tau}$ has the same intersections with  $\Phi(\sigma_{i_2})\circ \Phi(\sigma_{i_1})(\widetilde{\Delta})$ and with $\Phi(\sigma_{i_2}\sigma_{i_1})(\widetilde{\Delta})$, it follows that the two images of $\widetilde{\Delta}$ are istotopic on $\ws_{\sigma_{i_1}\sigma_{i_2}\cdot \tau}$.  Hence, Equation~\ref{eq} holds.



\end{proof}

\end{document}
